% Style: terse technical prose, British English (-ise). Code identifiers,
% SASS/PTX mnemonics, and named APIs stay US English ASCII verbatim.
% Author-style notes (from Connolly XP2010, ECBS2008): cold-open with a
% load-bearing fact, no throat-clearing; cite at the claim, quote borrowed
% phrases inline; impersonal voice (no "we" inflation); name tools plainly,
% no marketing adjectives; hedge honestly in-line; future work is specific
% or absent; short sections that earn their headers.
\documentclass[11pt]{article}

\usepackage[margin=1in]{geometry}
\usepackage{booktabs}
\usepackage{graphicx}
\usepackage{url}
\usepackage[numbers]{natbib}

\title{An Instruction-Level Characterisation of the NVIDIA TU102 GPU\\
\large Latency, Throughput, Memory Hierarchy, and NVLink Interconnect on sm\_75}
\author{David Connolly\\\texttt{david@connol.ly}}
\date{Draft — \today}

\begin{document}
\maketitle

\begin{abstract}
% Drafted at M6. One paragraph: Agner-style op table for TU102 (sm_75),
% methodology (dependent chains, fine-grained P-chase, SASS-gated benches,
% locked clocks), priors from Jia et al. T4 with per-row deviation tracking,
% additive NVLink/PCIe/NCCL characterization, and worked examples driving
% real kernel decisions in a production inference stack.
\end{abstract}

\section{Introduction}
\label{sec:introduction}

A used Quadro RTX 6000 sells for roughly a tenth of its launch price, and
two of them bridged by NVLink put 48~GB of GPU memory and a measured
1.2~TB/s of aggregate DRAM bandwidth on a desk. Whether that hardware can
carry a modern inference workload is not answerable from marketing
geometry: it depends on instruction latencies, cache behaviour, and
interconnect costs that NVIDIA does not publish and that existing
microbenchmark studies cover only for neighbouring chips. Jia et
al.\ dissected the TU104-based Tesla T4 \citep{jia2019turing}; the larger
TU102, with its different L2, DRAM system, and NVLink, has no equivalent
public characterisation.

This paper measures one. In the style of Agner Fog's x86 instruction
tables \citep{fog2025instruction}, it builds a complete operation table
for the TU102 at compute capability 7.5 (instruction latency and
throughput with measured issue-pipe bindings, the full memory hierarchy
from the register file to DRAM, and the NVLink, PCIe, and NCCL
interconnect costs) from microbenchmarks whose every timed loop is
verified at the SASS level and whose every published row carries its
provenance.

The research question is operational rather than encyclopaedic:
\emph{can a table of independently measured primitive costs predict the
behaviour of real composite kernels well enough to drive engineering
decisions?} The answer, developed through three pre-registered hypotheses
(\cref{sec:hypotheses}), is sharply split. Composed predictions are
excellent where a single resource binds, within 3.5\% on
single-resource kernels and within 10\% on a latency-dominated
interconnect primitive, and fail systematically where work couples across issue
pipes or where a primitive row turns out to be a floor rather than a
typical cost. Both failure modes are measured, diagnosed, and published.

The contributions are: (i) the table itself, 264 rows with per-row
provenance, priors, and deviation flags, regenerable from a fresh clone;
(ii) the measurement methodology, including a SASS-gated benchmark
discipline that caught every compiler-induced measurement defect before a
number shipped, and a clock-domain contract for a GPU whose memory clock
cannot be locked; (iii) findings absent from public documentation, among
them the fma-pipe binding of \texttt{IDP.4A}, the full-rate
f32-accumulate of Quadro-positioned TU102 tensor cores, the L1-retention
semantics of \texttt{.cs} loads, and the request-path saturation of the
L2; and (iv) a worked demonstration that pre-registration (gates,
metrics, and decision rules committed to a public repository before data
collection) transfers from empirical software engineering to
microbenchmarking, where the temptation to tune until numbers match
priors is the central threat to validity.

\section{Related work}
\label{sec:related}

Agner Fog's instruction tables \citep{fog2025instruction} are the model
for the artefact: exhaustive per-instruction latency and throughput for
x86 microarchitectures, maintained as data rather than narrative, with
the measurement method documented separately \citep{fog2025microarch}.
This work transplants the form to a GPU and adds two things the x86
tables do not carry: per-row provenance binding each number to the
benchmark commit and results line that produced it, and an explicit
deviation column against published priors.

GPU microbenchmarking begins for our purposes with Wong et
al.\ \citep{wong2010demystifying}, who established pointer-chase latency
measurement on GT200, and Mei and Chu \citep{mei2017dissecting}, whose
fine-grained pointer chase resolved cache geometry on Kepler and Maxwell.
Jia et al.\ characterised Volta \citep{jia2018volta} and then Turing via
the Tesla T4 \citep{jia2019turing}. Their instruction-latency tables are
the priors this paper compares against, row by row, where the prior
applies. The applicability boundary matters: the T4's TU104 shares the
TU102's streaming multiprocessor, so core-domain priors transfer, but its
L2, DRAM system, and lack of NVLink make memory-system and interconnect
priors non-binding, a distinction this paper records per row rather
than per paper.

Two methodological differences from Jia et al.\ are deliberate. First,
their measurements were taken with a custom SASS assembler, giving direct
control of the instruction stream. This work stays within the public
toolchain (PTX inline assembly under \texttt{nvcc} 13.3) and instead
\emph{verifies} the emitted stream, gating every benchmark on a
disassembly check. The compiler fights this approach: constant
folding, uniform-datapath conversion, strength reduction, and
approximate-arithmetic algebra each deleted or rewrote measurement
kernels during development. \Cref{sec:methodology} catalogues the
countermeasures, because any future toolchain-bound measurement effort
will face the same opponents. Second, where their report states results,
this one also pre-registers the decision-grade questions
(\cref{sec:hypotheses}) and publishes the failures.

\section{Methodology}
\label{sec:methodology}

All numbers come from one rig: a Dell T5820 (Xeon W-2140B) holding two
Quadro RTX 6000 GPUs, both at PCIe 3.0~x16, bridged by a two-link NV2
NVLink, driver 610.43.02, CUDA toolkit 13.3.33, ECC disabled on both
devices, CPU governor pinned to \texttt{performance}. The table is a
snapshot of this toolchain and driver pairing, not a living document. The
append-only results layout admits later datasets without disturbing this
one.

\subsection{Clock contract}
\label{sec:clocks}

Cycle-domain rows require a fixed SM clock. The harness prechecks a
1455~MHz lock (\texttt{nvidia-smi -lgc}) at startup and exits otherwise.
The memory clock cannot be locked on this part: \texttt{-lmc} reports
unsupported and the applications-clock interface is deprecated in this
driver. What holds instead, measured on both GPUs, is that CUDA compute
always executes in the P2 performance state with the memory clock at
6500~MHz. The harness ramps into P2 with a warmup kernel, verifies the
memory clock under load, and samples both clocks around every repetition,
rejecting any where either is off target. Bandwidth expectations are
therefore stated against the P2 memory clock: a 624~GB/s DRAM ceiling
rather than the 672 implied by the unreachable P0 point.

\subsection{Timing}

Latency rows use \texttt{clock64()} around dependent chains, with loop
overhead cancelled by a doubled-trip-count slope and the loop body
verified by size to fit the L0 instruction cache. Pointer-chase rows
embed full 64-bit pointers (or byte offsets, in shared and constant
windows) so the timed step is a single dependent load with no address
arithmetic \citep{wong2010demystifying,mei2017dissecting}. Throughput
rows are also timed in actual SM cycles, by one block on one SM:
wall-clock event timing normalised by the nominal locked clock showed a
0.3--1.0\% between-run spread that grew with region length (real-clock
thermal sag under the lock), where on-device cycle counts are immune.
Host-domain rows (launches, events, NCCL calls) use
\texttt{steady\_clock} over a thousand repetitions with a device-event
cross-check, reported as median with the 10th and 90th percentiles.

\subsection{SASS verification}
\label{sec:sass}

Operation selection is pinned by PTX inline assembly, and
\texttt{check\_sass.py} disassembles every benchmark binary, locates each
timed loop, and enforces that it contains the intended instructions and
nothing unexplained. The gate is not a formality. Working through the
table, \texttt{ptxas} or the NVVM front end variously: constant-folded
integer chains built on literal operands; moved warp-uniform integer
chains onto the uniform datapath (\texttt{UPOPC} on \texttt{UR}
registers); strength-reduced dependent add chains into \texttt{IMAD} and
\texttt{LEA} through both a dead carry-out pin and cross-coupled
accumulators; deleted six of eight independent accumulator chains whose
results a sink did not consume; folded \texttt{rcp.approx} compositions
as exact algebra, twice; contracted a deliberately separate add into an
\texttt{FFMA}; lowered an f16 widening conversion to \texttt{HADD2}
rather than the expected \texttt{F2F}; and expanded a 64-bit remainder
into a division subroutine call whose float-reciprocal emulation put
\texttt{MUFU} and conversion traffic inside a bandwidth loop. Every one
of these was caught by counting emitted mnemonics against the design,
not by a timing looking wrong. Several produced plausible numbers for
the wrong instruction stream. Both gates carry negative tests: a
deliberately unlocked clock and a deliberately miscompiled kernel must
fail.

\subsection{Statistics and gate taxonomy}

Each row is the median of ten repetitions, with the coefficient of
variation published. Every published row further requires at least two
independent process invocations on each GPU, and the between-run spread
must not exceed the within-run variation beyond a domain floor (0.1\% for
cycle rows, 0.5\% for wall-clock bandwidths, 5\% for host-domain times).
SM-domain rows must agree across the two physical GPUs within their
combined variation. Interconnect rows differ by design and carry
direction in the variant. Verification gates are of two kinds, and
confusing them is how benchmark suites quietly tune themselves to the
literature: \emph{methodology-sanity} gates (two independent methods must
agree) block publication on failure, while \emph{prior-comparison} gates
(against \citealp{jia2019turing} or whitepaper geometry) publish the
deviation whatever it is. The table exists to report true deviations.

\subsection{Pre-registration}
\label{sec:prereg}

The coverage manifest and measurement policy were committed to the public
repository before any benchmark ran, and the three decision-grade
hypotheses of \cref{sec:hypotheses}, with their predictions,
operationalisations, and decision rules, were committed before their
data existed (repository commit \texttt{a4fdd73}, 2026-06-10, preceding
every measurement it governs). The commit history is the receipt. One
registered rule fired mid-study and was honoured: a composed-prediction
gate failed at its registered tolerance, and the affected comparison was
blocked rather than re-derived, until independently measured constituents
allowed a documented remediation (\cref{sec:hypotheses-outcomes}).

\subsection{A worked row}

The discipline end to end, for \texttt{alu.ffma.lat}. A 128-deep
dependent \texttt{fma.rn.f32} chain on runtime operands compiles to
exactly 128 \texttt{FFMA} in the timed loop, and the gate passes. Ten
repetitions at two trip counts give a slope of 4.07 cycles with zero
within-run variation, and four invocations across both GPUs agree within
the 0.1\% floor. Against the 4-cycle prior
\citep[tab.~4.1]{jia2019turing} the row publishes as 4.07, deviation
$+1.8\%$, flag \texttt{ok}, with provenance naming the benchmark file,
commit, and results lines.

\section{SM Pipelines}

% Drafted at the milestone that produces this section's data.

\section{Memory Hierarchy}

% Drafted at the milestone that produces this section's data.

\section{Interconnect}

% Drafted once the data for this section has been measured.

\section{Worked Examples}

% Drafted once the data for this section has been measured.

\section{Conclusion}

% Drafted at the milestone that produces this section's data.

\section*{Acknowledgments}

The microbenchmark harness, analysis tooling, and manuscript drafts were
developed with the assistance of Claude (model \texttt{claude-fable-5},
Anthropic), used via Claude Code. All measurements were executed, validated,
and interpreted on the author's hardware under the author's direction; the
author takes sole responsibility for the content.

The author claims no affiliation for this work.


\bibliographystyle{plainnat}
\bibliography{references}

\end{document}
